%% kapa: exam=1 round=35 year=2024 period=1 form=A subject=civil_law subject_ko=민법 seq=36 number=36
%% source: Q-Net 감정평가사 기출문제 (cst00309) artlSeq=5213229 file=raw/1008/5213229/1교시 문제지 원본.pdf
%% conversion: kordoc 4.24.51 -> tools/gen1.py (auto); review=proofread-sonnet; flags=none
\begin{question}{36}
유치권이 유효하게 성립할 수 있는 경우는? (다툼이 있으면 판례에 따름)
\begin{choices}
  \item 주택수선공사를 한 수급인이 공사대금채권을 담보하기 위하여 주택을 점유한 경우
  \item 임대인이 지급하기로 약정한 권리금의 반환청구권을 담보하기 위하여 임차인이 상가 건물을 점유한 경우
  \item 매도인이 매수인에 대한 매매대금채권을 담보하기 위하여 매매목적물을 점유한 경우
  \item 주택신축을 위하여 수급인에게 공급한 건축자재에 대한 대금채권을 담보하기 위하여 그 공급자가 주택을 점유한 경우
  \item 임차인이 임차보증금반환채권을 담보하기 위하여 임차목적물을 점유한 경우
\end{choices}
\end{question}
\answer{1}
